\begin{problem}[第31届全国中学生物理竞赛复赛][等腰三角形平板平衡]{19}
如图所示，一质量为 $m$、底边 AB 长为 $b$、等腰边长为 $a$、质量均匀分布的等腰三角形平板，可绕过光滑铰链支点 A 和 B 的水平轴 $x$ 自由转动；图中原点 O 位于 AB 的中点，y 轴垂直于板面斜向上，$z$ 轴在板面上从原点 O 指向三角形顶点 C。今在平板上任一给定点 $\mathrm{M}_0(x_0,0,z_0)$ 加一垂直于板面的拉力 $Q$。
    \begin{subproblem}
        若平衡时平板与竖直方向成的角度为 $\varphi$，求拉力 Q 以及铰链支点对三角形板的作用力 $N_A$ 和 $N_B$；
    \end{subproblem}
    \begin{subproblem}
        若在三角形平板上缓慢改变拉力 $Q$ 的作用点 M 的位置，使平衡时平板与竖直方向成的角度仍保持为 $\varphi$，则改变的作用点 M 形成的轨迹满足什么条件时，可使铰链支点 A 或 B 对板作用力的垂直平板的分量在 M 变动中保持不变？
    \end{subproblem}
\end{problem}